\documentclass[10pt,a4paper]{amsart}
\usepackage[margin=19mm]{geometry}
\usepackage{amsmath,amssymb,mathtools}
\usepackage[scr=rsfs,cal=euler]{mathalfa}
\usepackage{xfrac}
\usepackage[unicode,hidelinks,bookmarks=false]{hyperref}
\newcommand{\ie}{i.e.\ }
\newcommand{\cf}{cf.\ }
\newcommand{\resp}{respectively\ }
\newcommand{\Subsection}{\subsection}
\providecommand{\nobreakdash}{\nobreak\hbox{-}}
\setlength{\emergencystretch}{2em}
\multlinegap0pt
\usepackage{amssymb}
\newtheorem{theorem}{Theorem}
\def\P{{{\mathbb P}}}
\def\Q{{{\mathbb Q}}}
\title{\uppercase{Rational and integral values of rational functions at rational points}}
\author{Pietro Corvaja, Umberto Zannier}
\date{Source: arXiv:2608.28255v1 (CC BY 4.0)}
\begin{document}
\maketitle

\noindent\emph{Source and licence.} \uppercase{Rational and integral values of rational functions at rational points}. Version arXiv:2608.28255v1 (CC BY 4.0), distributed by arXiv under the Creative Commons Attribution 4.0 International licence, http://creativecommons.org/licenses/by/4.0/, as recorded in the arXiv licence metadata for this exact version and shown on the abstract page https://arxiv.org/abs/2608.28255v1. Full licence notice: \texttt{licenses/arxiv-2608.28255-CC-BY-4.0.txt}.\par\smallskip

\noindent\textit{Editorial scope.} Counted unit: the article's theorem T.Main2 (source0.tex lines 877--879). The article's complete original proof of it is delivered with it (source0.tex lines 905--928, one \texttt{proof} environment of 24 lines, which the article places away from the statement). No mathematics is written, repaired or re-derived: the statement and the proof are the article's own lines, in the article's own notation, and no retained line is substituted or re-flowed.  No further source-local context is needed: every cross-reference of every retained line resolves inside the two delivered spans.  Nothing was omitted from any retained span, so the package carries no omission record.  No retained line contains a citation, so the wrapper carries no bibliography and no bibliography entry.  No retained line contains a figure environment, an image-inclusion command or drawing code, so no image is delivered and the package declares no asset dependency.  Editorial formatting only, in addition: an AMS \texttt{amsart} preamble that restates the article's own declarations and macros used by the retained lines, namely the environments \texttt{theorem} on separate counters, together with the standard packages those lines need.  The retained environments print in this document as Theorem 1 in order of appearance, whereas the article prints the same items with the numbering of its own full document.  The complete native source of the article as distributed in the arXiv e-print for this exact version is bundled unchanged as \texttt{sources/208775/source0.tex}.
\begin{theorem}\label{T.Main2}
Let $X$ be a a semi-abelian variety defined over a number field $k$; let $f:X\to \P_1$ be a non-constant map. Let $\Gamma\subset X(k)$ be a finitely generated subgroup. The set of points in $X(k)$ which are not on the image $f(\Gamma)$ satisfy the weak-approximation property. In particular, they form an infinite set.  
\end{theorem}

\begin{proof}[ Proof of Theorem \ref{T.Main2}] 

Let $X,k,f$ be as in the statement. Let us choose a finite set of non-archimedean places $\{\nu_1,\ldots,\nu_h\}$ of $k$, rational points $q_1,\ldots,q_k$ and positive integers  $r_1,\ldots,r_h$. We look for points $p \in \P_1(k)\setminus f(\Gamma)$ which  approximate   the points $p_1,\ldots,p_h$ with respect to the given valuations, i.e. satisfy
\begin{equation}\label{E.approximation}
p\equiv q_i \pmod {\nu_i^{r_i}} \qquad \mathrm{for} \ i=1,\ldots,h.
\end{equation} 
By the weak approximation property on the line (i.e. the Chinese remainder theorem), we can find a point $p_0\in \P_1(k)$ such that the above system of congruences is satisfied for $p=p_0$.

Then, construct an elliptic curve $E$ over $k$, of positive Mordell-Weil rank, and a morphism $\pi:E\to \P_1$, still defined over $k$, such that $\pi$ sends the neutral element of $E$ to $p_0$. We also require that $E$ admits no non-constant morphism to $X$. 
It is clear that such an elliptic curve can be constructed: by Poincar\'e's complete reducibility theorem   the elliptic curves admitting non-constant morphisms to $X$ lie in only finitely many isogeny classes and then by the previous lemma one can pick an elliptic curve over $k$ (even over $\Q$) outside this finite union of classes, having positive Mordell-Weil rank. 

Now, we claim that only finitely rational points of $\P_1$ can lie both in the value sets   $\pi(E(k))$ and $f(\Gamma)$: indeed, consider the fibre product 
$$
Y=X\times_{f,\pi} E=\{(x,u)\in X\times E\, :\, f(x)=\pi(u)\}\subset X\times E.
$$
It is a hypersurface of the semi-abelian variety $X\times E$ and is endowed with a natural map to $\P_1$, sending $(x,u)\mapsto f(x)=\pi(u)$. The points in $f(X(k))\cap \pi(E(k))$ are precisely the images of the rational points in $Y$. Now, by Faltings-Vojta's theorem on rational points on sub-varieties of semi-abelian varieties, the set $Y(k)$ is contained in the union of finitely many translates of algebraic subgroups of $X\times E$ contained in $Y$.

By our assumption on $E$ the only connected algebraic subgroups of the product $X\times E$ are of the form  $A\times E$ or $A\times\{0_E\}$,  where $A\subset X$ is an algebraic subgroup of $X$. Let us show that  $Y$ does not contain any translate of a  subgroup of the form $A\times E$. Indeed,   if for some $x_0 \in X\times E$ the translate $(u_0+A)\times E$  would be contained in $Y$, from the equation $f(u_0+a)=\pi(u)$ valid for all $a\in A$ and $u\in E$, we would obtain a contradiction (since $\pi$ is non-constant).

Then all rational points in $Y(k)$ are contained in sets of the form $X(k)\times F$, for a finite set $F\subset E(k)$. In other words, the projection $\pi: X\times E \to \P_1$ takes only finitely many values in the points of $Y(k)$. Equivalently, for all but finitely many points $u\in E(k)$, the rational point $\pi(u)$ does not lie in the value set $f(X(k))$, and {\it a fortiori} not in the image $f(\Gamma)$. 

To construct the rational point $p\in \P_1(k)\setminus f(\Gamma)$ approximating the given points $q_1,\ldots,q_h$ as in \eqref{E.approximation}, observe that $p=\pi(0_E)$ does solve the system \eqref{E.approximation}. Take now any non-torsion point $u_0\in E(k)$; its multiples $nu_0$, for sufficiently divisible $n\in Z$, are congruent to the neutral element $0_E$ modulo $\nu_i^{r_i}$ for $i=1,\ldots,h$. Then their images $\pi(nu_0)$ solve the congruence system \eqref{E.approximation}, and all but finitely many of them lie outside the set $f(\Gamma)$.

\end{proof}

\end{document}
